\documentclass[12pt, a4paper]{article}
\usepackage{ctex}
\usepackage{geometry}
\usepackage{amsmath, amssymb}
\usepackage{graphicx}
\usepackage{booktabs}
\usepackage{hyperref}
\usepackage{listings}
\usepackage{xcolor}
\usepackage{fancyhdr}

\geometry{top=2.5cm, bottom=2.5cm, left=2.5cm, right=2.5cm}

% 页码设置
\pagestyle{fancy}
\fancyhf{}
\cfoot{\thepage}
\renewcommand{\headrulewidth}{0pt}

% 代码块样式
\lstset{
    basicstyle=\small\ttfamily,
    breaklines=true,
    frame=single,
    language=Python
}

\begin{document}

% ===== 第一页：承诺书（从官网下载后手动插入）=====
\newpage

% ===== 第二页：编号专用页（从官网下载后手动插入）=====
\newpage

% ===== 第三页：摘要页 =====
\setcounter{page}{1}
\begin{center}
    {\Large \textbf{论文标题}}\\[0.5cm]
    \textbf{摘要}
\end{center}

本文针对……问题，建立了……模型。
对于问题一，……；
对于问题二，……；
对于问题三，……。
最终得到……结论。

\vspace{0.5cm}
\noindent\textbf{关键词：} 关键词1；关键词2；关键词3

\newpage

% ===== 正文 =====

\section{问题重述}
% 用自己的话概括题目，明确各子问题

\section{问题分析}
% 分析题目类型，说明整体解题思路

\section{模型假设}
\begin{enumerate}
    \item 假设一……，由：……
    \item 假设二……，由：……
    \item 假设三……，由：……
\end{enumerate}

\section{符号说明}
\begin{table}[h]
\centering
\begin{tabular}{cl}
\toprule
符号 & 说明 \\
\midrule
$x$ & 说明 \\
$y$ & 说明 \\
\bottomrule
\end{tabular}
\end{table}

\section{模型建立与求解}

\subsection{问题一：XXX}

\subsubsection{建模思路}

\subsubsection{模型建立}

\subsubsection{求解结果}

\subsection{问题二：XXX}

\subsubsection{建模思路}

\subsubsection{模型建立}

\subsubsection{求解结果}

\subsection{问题三：XXX}

\subsubsection{建模思路}

\subsubsection{模型建立}

\subsubsection{求解结果}

\section{模型评价与改进}

\subsection{模型优点}
\begin{itemize}
    \item 
\end{itemize}

\subsection{模型局限性}
\begin{itemize}
    \item 
\end{itemize}

\subsection{改进方向}

\begin{thebibliography}{99}
\bibitem{ref1} 作者. 标题[J]. 期刊, 年份.
\end{thebibliography}

% ===== 附录 =====
\appendix
\section{支撑材料文件列表}
\begin{itemize}
    \item \texttt{code/} — 全部源代码
    \item \texttt{data/} — 数据文件
\end{itemize}

\section{核心代码}
\begin{lstlisting}
# 在此粘贴核心代码
\end{lstlisting}

\end{document}